\documentclass[12pt, titlepage]{article}
%\usepackage{hyperref}

\title{Feedback Controls Lab Project Proposal \\ \normalsize{EGEE 4810 - Senior Design}}
\author{Aydan Vissing, Brady Schick, Gabriel Southwell, Simon Ross}
\date{\footnotesize{Due} \\ \normalsize{September 12, 2025}}


%\setlength{\topmargin}{-.25in}
%\setlength{\oddsidemargin}{.25in}
%\setlength{\evensidemargin}{.25in}
%\setlength{\textwidth}{6.0in}
%\setlength{\textheight}{9.0in}
%\renewcommand{\baselinestretch}{1.4}
\usepackage{indentfirst} %Maybe keep this maybe not

\begin{document}
\maketitle

%FIXME remove comments

%FIXME improve abstract
\begin{abstract}
Lorem Ipsum
\end{abstract}

\newpage
\tableofcontents
\listoffigures
\listoftables

\newpage
\section{Definition Statement}
1-2 sentence description of why we are here
\section{Background}
The current lab used for EGME-4660 uses the Feedback Analogue Unit 33-110 and the Feedback Mechanical Unit 33-100. Although the labs using the equipment teach the material necessary, they are often tedious. The user interface to conduct the experiments is a circuit board with several banana plug chords, which may seem boring and irrelevant to the mechanical engineering students. In addition, adjusting the power supply or wiring the power incorrectly can result in damage to the equipment, often requiring parts to be replaced. We plan to improve the lab experience by replacing the current equipment with drones that can be reprogrammed for the purposes of the lab with a friendly user interface.
\section{Objectives}
\subsection{Marketing Requirements}
\begin{itemize}
  \item[$\Rightarrow$] Lab Experience Must:
  \begin{itemize}
    \item Focus on autonomous actuation and control
    \item Include a PID control design experience
    \item Visually and graphically display a step response
    \item Measure a Bode plot in real life
    \item Include physical and observable disturbances/noise
    \item Showcase the results physically in real life (not just with numbers/data)
    \item Ability to log data to a neutral format independent of real time visualizations
    \item Include widgit to generate clean plots of the data
    \item Include capability for a tracking experiment
    \item Include relevant safety measures
  \end{itemize}
\end{itemize}
\subsection{Engineering Requirements}
\subsubsection{Hardware Requirements}
\begin{itemize}
  \item[$\Rightarrow$] Hard Requirements:
  \begin{itemize}
    \item Microcontroller
    \item 4 motors
    \item Battery power supply
    \item PCB
    \item Bladeguards for safety
  \end{itemize}
  \item[$\Rightarrow$] Soft Requirements:
  \begin{itemize}
    \item Sturdy frame
    \item Conveniently swappable and rechargeable batteries
  \end{itemize}
\end{itemize}
\subsubsection{Software Requirements}
\begin{itemize}
  \item[$\Rightarrow$] Hard Requirements:
  \begin{itemize}
    \item PID controller logic so drone can fly
    \item Memory for data collection
    \item User-friendly GUI for programming
    \item Software to plot collected data
  \end{itemize}
  \item[$\Rightarrow$] Soft Requirements:
  \begin{itemize}
    \item Wireless communication
  \end{itemize}
\end{itemize}
\subsection{Constraints}
\begin{itemize}
  \item Labs conducted in ENS 145
  \item \$600 budget
  \item Lab safety
  \item No MATLAB
\end{itemize}

\section{Officer Designations}
To aid in the organization of this project, we have outlined four officer designations below (one for each member of the team). The goal of this is to delegate various essential tasks to individuals so that they are not overlooked and each member is not overloaded. The roles will likely be rather fluid since the team is small, but give general guidelines on where to look for work. These roles are as follows: \\

\textbf{CFO - Aydan Vissing} The Chief Financial Officer (CFO) role is in charge of managing finances for the project, heading up budgets, and making sure the project stays within the budget. This role will be deeply involved in the use of time and resources throughout the project. \\

\textbf{CTO - Brady Schick} The Chief Technical Officer (CTO) role will be responsible for managing the hardware side of the project, with an emphasis on keeping the big picture in hardware in perspective. This role will likely shift to a focus in software as the project matures. \\

\textbf{CSO - Gabriel Southwell} The Chief Software Officer (CSO) will make sure that software development proceeds according to plan, ensuring that proper hardware is acquired for the software needs, software is organized and able to be applied to the microcontroller, and UI developments are user-friendly and intuitive. This role will likely increase in difficulty throughout the course of the project. \\

\textbf{CEO - Simon Ross} The Chief Executive Officer (CEO) is responsible for paperwork, organization, and keeping the big picture of the project in perspective. This role will make sure that the members act as a unified group and that timely progress is made towards the project milestones. \\

Again, these roles simply provide a basic structure for the team so that essential duties are not forgotten - each member will likely have a large amount of overlap in their duties in the project.

\section{Strategy}
Hardware and Software

\section{Project Schedule}

\section{Milestones}

\section{Reporting Plan}
Every Monday, weekly reports will be submitted with statements from each member of the team. This will keep a record of what has been accomplished each week and what will plan to be accomplished in the future. Reports will include both group and individual accomplishments.
  
Every Friday at 11:00 AM, weekly team meeting will be held with the project advisor/customer, Dr. Fredette. These meeting ensure that the project is continuing to progress in the desired direction.
  
On December 12th, a design report will be submitted to document the progress that has been made up to that point.
  
At the end of April, a final report will be submitted to document the entirety of the project.

\section{Testing Plan}
What things do we want to be able to test

\section{Budget}


\end{document}
